\documentclass[11pt]{article}
\usepackage[margin=1in]{geometry}
\usepackage{amsmath,amssymb,amsthm,graphicx,microtype,hyperref}
\hypersetup{colorlinks=true,linkcolor=blue,urlcolor=blue,citecolor=blue}

\newtheorem{theorem}{Theorem}
\newtheorem{lemma}{Lemma}
\newcommand{\R}{\mathbb R}
\newcommand{\Snn}{\mathbb S_{\mathrm{nn}}}
\newcommand{\cN}{\mathcal N}

\title{Polynomial Metric Entropy of the Monotone Spherical Chamber}
\author{Solution packet for arXiv:1210.0139}
\date{22 July 2026}

\begin{document}
\maketitle

\noindent\fbox{\parbox{0.92\textwidth}{\textbf{Status.}
Partial result likely valid, subject to human review.  The theorem determines
the sharp order in the ambient dimension at every fixed chordal accuracy
$0<\varepsilon\leq 1/5$: the optimal covering number is polynomial in $M$,
whereas the source construction is quasi-polynomial.}}

\section{Source question}

Fickus, Jasper, Mixon, and Peterson define
\[
 \Snn^{M-1}=\{x\in\R^M:0\leq x_1\leq\cdots\leq x_M,
 \ \|x\|_2=1\}
\]
and use the chordal distance
\[
 \rho(x,y)=\sqrt{1-|\langle x,y\rangle|^2}.
\]
Their signed-permutation reduction makes an efficient net for this chamber the
computational bottleneck in their NERF estimates.  In the conclusions they ask,
in effect, how many nonnegative monotone unit vectors are needed to approximate
all such vectors in chordal distance \cite[p.~16]{FJMP}.

\begin{figure}[ht]
\centering
\includegraphics[width=\textwidth]{figures/open_problem_crop.png}
\caption{The open-ended covering question in the source paper, page 16.}
\end{figure}

Let $\cN_M(\varepsilon)$ denote the least cardinality of an
$\varepsilon$-net for $\Snn^{M-1}$ in $\rho$.  The source constructs, for fixed
$\varepsilon$, a net of size
$\exp(O_\varepsilon((\log M)^2))$.  The result below improves the dimension
dependence to a power of $M$, and shows that a power is necessary at sufficiently
fine fixed accuracy.

\section{Proof intuition}

The upper bound uses monotonicity in the coordinate index rather than
quantization in the value axis.  Reverse the coordinates and divide them into
geometrically growing consecutive blocks.  On a block beginning at $a$, use at
most $\eta a$ coordinates.  Replacing a decreasing vector by its value at the
right endpoint of each block loses at most an $\eta$ fraction of its squared
mass.  Thus every chamber point is close to a fixed block-constant subspace of
dimension $O(\eta^{-1}\log M)$.  A standard net in that subspace gives the
upper bound after taking $\eta\asymp\varepsilon^2$.

For the lower bound, reverse the idea.  Use $d\asymp\log M$ blocks whose sizes
decrease by a factor $16$.  Give every block nearly equal \emph{energy}, and
perturb those energies by balanced signs from a binary code.  The factor $16$
leaves enough separation between adjacent block heights to preserve
monotonicity under every perturbation.  Codewords far apart in Hamming distance
then give chamber points a fixed positive chordal distance apart.

\section{The metric-entropy theorem}

\begin{theorem}\label{thm:main}
There are universal constants $c>0$ and $M_0\geq2$ such that, for every
$M\geq M_0$ and $0<\varepsilon\leq1/5$,
\begin{equation}\label{eq:main}
 M^c\leq \cN_M(\varepsilon)
 \leq
 \left(1+\frac4\varepsilon\right)^{
 32\varepsilon^{-2}(1+\log M)}.
\end{equation}
Consequently, for each fixed $0<\varepsilon\leq1/5$,
\[
 \log \cN_M(\varepsilon)=\Theta_\varepsilon(\log M)
 \qquad(M\to\infty).
\]
The upper bound remains valid for every $0<\varepsilon<1$.
\end{theorem}

Reversing the order of the coordinates is an isometry, so in the upper-bound
proof we work with
\[
 T_M=\{x\in\R^M:x_1\geq\cdots\geq x_M\geq0,\ \|x\|_2=1\}.
\]
For points in either chamber the inner product is nonnegative, and
\begin{equation}\label{eq:metriccomparison}
 \rho(x,y)^2=(1-\langle x,y\rangle)(1+\langle x,y\rangle)
 \leq 2(1-\langle x,y\rangle)=\|x-y\|_2^2.
\end{equation}

\begin{lemma}[geometric-block compression]\label{lem:compression}
For every $0<\eta\leq1$ there is a partition of
$\{1,\ldots,M\}$ into $d$ consecutive blocks such that
\begin{equation}\label{eq:dimbound}
 d\leq4\eta^{-1}(1+\log M),
\end{equation}
and the orthogonal projection $P$ onto the subspace of vectors constant on
each block satisfies
\begin{equation}\label{eq:projectionerror}
 \|x-Px\|_2^2\leq\eta\qquad(x\in T_M).
\end{equation}
Moreover, $Px$ is nonnegative and nonincreasing.
\end{lemma}

\begin{proof}
Make every index $a<2/\eta$ a singleton.  Thereafter, whenever the next block
starts at $a$, give it length
\[
 \ell_a=\min\{\lfloor\eta a\rfloor,M-a+1\}.
\]
For a nonterminal nonsingleton block,
$\eta a/2\leq\ell_a\leq\eta a$.  Its successor therefore begins at least at
$(1+\eta/2)a$.  Since
$\log(1+\eta/2)\geq\eta/4$, counting the initial singletons and subsequent
geometric growth gives \eqref{eq:dimbound}.

Fix $x\in T_M$.  On a block $B=[a,b]$, define $q_i=x_b$ for $i\in B$.
Then $q$ is block-constant, and
\begin{equation}\label{eq:squareloss}
 \sum_{i\in B}(x_i-x_b)^2
 \leq\sum_{i\in B}(x_i^2-x_b^2).
\end{equation}
Put $z_i=x_i^2$.  For any $t\geq0$, the set
$\{i:z_i>t\}$ is an initial interval $\{1,\ldots,k(t)\}$.  In the layer-cake
representation of the right side of \eqref{eq:squareloss}, complete blocks
cancel.  If $k(t)$ lies in a nonsingleton block $[a,b]$, the uncancelled number
of indices is at most $|B|\leq\eta a\leq\eta k(t)$; it is zero in an initial
singleton.  Hence
\[
 \|x-q\|_2^2
 \leq \int_0^\infty \eta k(t)\,dt
 =\eta\sum_{i=1}^M z_i=\eta.
\]
Orthogonal projection is at least as accurate as $q$, proving
\eqref{eq:projectionerror}.  Finally, $Px$ consists of the consecutive block
averages of $x$, so it remains nonnegative and nonincreasing.
\end{proof}

\begin{proof}[Proof of the upper bound in Theorem~\ref{thm:main}]
Take $\eta=\varepsilon^2/8$ and the subspace $V$ from
Lemma~\ref{lem:compression}.  If $e=\|x-Px\|_2$, orthogonality gives
$\|Px\|_2=\sqrt{1-e^2}$.  Thus, for
$u=Px/\|Px\|_2\in T_M\cap V$,
\begin{equation}\label{eq:normalization}
 \|x-u\|_2^2
 =2(1-\|Px\|_2)
 \leq2e^2\leq2\eta=\varepsilon^2/4.
\end{equation}

Take a maximal Euclidean $\varepsilon/2$-separated subset $A$ of
$T_M\cap V$.  It is an $\varepsilon/2$-net of that set.  The usual disjoint-ball
argument in the $d$-dimensional space $V$ gives
\[
 |A|\leq(1+4/\varepsilon)^d.
\]
Equation \eqref{eq:normalization} and the triangle inequality show that $A$ is
an Euclidean $\varepsilon$-net of $T_M$, and
\eqref{eq:metriccomparison} makes it a chordal $\varepsilon$-net as well.
Since Lemma~\ref{lem:compression} gives
$d\leq32\varepsilon^{-2}(1+\log M)$, the upper estimate follows.
\end{proof}

\begin{proof}[Proof of the lower bound in Theorem~\ref{thm:main}]
Let $d$ be the largest even integer for which
\[
 m_d:=1+16+\cdots+16^{d-1}=\frac{16^d-1}{15}\leq M.
\]
Thus $d\geq c_1\log M$ for all sufficiently large $M$.  After $M-m_d$
initial zero coordinates, split the remaining coordinates into consecutive
blocks $B_1,\ldots,B_d$ of sizes
$|B_j|=16^{d-j}$.

We use a family $\mathcal C\subseteq\{-1,1\}^d$ of balanced words
($\sum_j\sigma_j=0$) with pairwise Hamming distance at least $d/4$ and
$|\mathcal C|\geq\exp(c_2d)$.  For completeness, such a family follows by a
greedy selection from the $\binom d{d/2}$ balanced words.  A Hamming ball of
radius less than $d/4$ contains at most
\[
 \sum_{r\leq d/8}\binom{d/2}{r}^2
 \leq(d+1)2^{dH_2(1/4)}
\]
balanced words, while
$\binom d{d/2}\geq2^d/(d+1)$ and $H_2(1/4)<1$.

For $\sigma\in\mathcal C$, define $x^\sigma$ to be constant on $B_j$ with
value
\begin{equation}\label{eq:packingvector}
 x^\sigma_i=
 \frac{1+\sigma_j/2}{\sqrt{(5/4)d\,|B_j|}}
 \qquad(i\in B_j).
\end{equation}
Balancedness gives $\|x^\sigma\|_2=1$.  Adjacent block heights obey
\[
 \frac{x^\sigma|_{B_{j+1}}}{x^\sigma|_{B_j}}
 =4\frac{1+\sigma_{j+1}/2}{1+\sigma_j/2}\geq\frac43,
\]
so every $x^\sigma$ belongs to the nonnegative nondecreasing chamber.

If $h$ is the Hamming distance between $\sigma$ and $\tau$, direct calculation
using balancedness gives
\[
 \langle x^\sigma,x^\tau\rangle
 =1-\frac{2(1/2)^2}{1+(1/2)^2}\frac hd
 =1-\frac25\frac hd.
\]
For $h\geq d/4$, this is at most $9/10$, and hence
\[
 \rho(x^\sigma,x^\tau)
 \geq\sqrt{1-(9/10)^2}=\frac{\sqrt{19}}{10}>\frac25.
\]
When $\varepsilon\leq1/5$, an $\varepsilon$-ball contains at most one member
of this packing.  Therefore
\[
 \cN_M(\varepsilon)\geq|\mathcal C|
 \geq\exp(c_2d)\geq M^c,
\]
after decreasing the universal constant $c>0$ and increasing $M_0$.
\end{proof}

\section{Verification and limitations}

The argument is elementary.  The accompanying script
\texttt{code/verify\_entropy\_bounds.py} checks the geometric-block estimate on
random and adversarial monotone vectors, verifies the dimension count, and
enumerates finite balanced-code packings to check normalization, monotonicity,
and the chordal separation constant.  These computations are sanity checks,
not ingredients of the proof.  It passed with 300 random trials at
$(M,\varepsilon,d)=(4096,0.2,8)$ and 200 trials at
$(7777,0.125,10)$.  The commands were
\begin{verbatim}
python code/verify_entropy_bounds.py --dimension 4096 --epsilon 0.2 \
  --trials 300 --packing-blocks 8
python code/verify_entropy_bounds.py --dimension 7777 --epsilon 0.125 \
  --trials 200 --packing-blocks 10
\end{verbatim}
In the first run, the minimum enumerated packing distance was exactly
$\sqrt{19}/10$, the constant used in the proof.

The theorem settles the dimension dependence of the optimal covering number at
every fixed $0<\varepsilon\leq1/5$.  It does not determine the sharp dependence
of the exponent on $\varepsilon$, and the maximal-net step is an existence
construction rather than a practically optimized enumeration.  Accordingly,
this packet is classified as a substantial partial result rather than a full
answer to every algorithmic interpretation of the source's broad question.

A bounded novelty search used the run indexes, arXiv:1210.0139, the exact source
phrases ``nonnegative, nonincreasing unit norm'' and ``more efficient
$\varepsilon$-nets,'' and combinations of ``metric entropy,'' ``monotone
cone,'' ``isotonic cone,'' and ``chordal distance.''  No later paper directly
settling this chamber covering question was found.  This is a bounded search,
not an exhaustive priority claim.

\paragraph{Human-review recommendation.}
Check the layer-cake cancellation in Lemma~\ref{lem:compression}, the
normalization estimate \eqref{eq:normalization}, and the constant-weight-code
packing.  These are the only substantive steps.

\begin{thebibliography}{9}
\bibitem{FJMP}
M. Fickus, J. Jasper, D. G. Mixon, and J. Peterson,
\emph{Group-theoretic constructions of erasure-robust frames},
Linear Algebra Appl. 479 (2015), 131--154; arXiv:1210.0139.
\end{thebibliography}

\end{document}
